\section{The frozen tensor-subset shift and its parity fiber}
\label{sec:source}

\subsection{Tensor source and signed alphabet}

Let
\[
\M=\{F_n:n\ge1\},\qquad
F_m\otimes F_n=F_{mn},\qquad h(F_n)=\log n.
\tag{3.1}\label{3.1}
\]
The tensor-indecomposable elements are \(F_p\).  Fix a finite atom set \(P\)
and assign a commuting variable \(x_p\) to each atom.  The alphabet and full
shift are
\[
\E_P=\{S:\varnothing\ne S\subseteq P\},\qquad
X_P=\E_P^{\Z},\qquad \sigma:X_P\to X_P.
\tag{3.2}\label{3.2}
\]
For \(S\in\E_P\), set
\[
x_S=\prod_{p\in S}x_p,\qquad
\eps(S)=(-1)^{|S|+1},\qquad
w(S)=\eps(S)x_S.
\tag{3.3}\label{3.3}
\]
The sign in \eqref{3.3} is an ordinary scalar edge coefficient.  On a temporal
repetition it is raised to the corresponding scalar power; it is not replaced
by a chain or supertrace convention.

Arithmetic specialization takes \(x_p=p^{-s}\) and
\[
T(S)=\sum_{p\in S}\log p,\qquad
w_s(S)=\eps(S)e^{-sT(S)}.
\tag{3.4}\label{3.4}
\]
The polynomial identities are established before this specialization.  Thus
primality is part of the source interpretation but not an ingredient in the
factorization proof.

\subsection{The intrinsic \texorpdfstring{\(C_2\)}{C2} extension}

Write \(C_2=\{0,1\}\) additively and define
\[
\alpha_P(S)=|S|\pmod2.
\tag{3.5}\label{3.5}
\]
The skew product is
\[
\sigma_\alpha(x,g)=\bigl(\sigma x,g+\alpha_P(x_0)\bigr).
\tag{3.6}\label{3.6}
\]
For \(h\in C_2\), deck translation
\(R_h(x,g)=(x,g+h)\) satisfies
\[
R_h\sigma_\alpha(x,g)
=\bigl(\sigma x,g+\alpha_P(x_0)+h\bigr)
=\sigma_\alpha R_h(x,g).
\tag{3.7}\label{3.7}
\]
Equation \eqref{3.7} remains true after every roof specialization because
\(R_h\) acts only on the fiber.  It never permutes \(P\), \(x_p\), or \(T(S)\).

\begin{proposition}[Transitivity and mixing]
\label{prop:mixing}
For every nonempty finite \(P\), the \(C_2\) extension is topologically
transitive.  It has period two when \(|P|=1\) and is mixing when \(|P|\ge2\).
\end{proposition}

\begin{proof}
Every singleton symbol has odd degree and switches the two fiber states, so
the two-state presentation is strongly connected.  If \(|P|=1\), every edge
switches the fiber and every closed path has even length.  If \(|P|\ge2\), an
even subset symbol gives a loop at each fiber state.  Strong connectivity
together with a length-one loop makes the period one.
\end{proof}

The singleton fixed point has cocycle product \(1\in C_2\), which is
nonidentity.  Since a vertex coboundary telescopes to identity on every
periodic orbit, the parity cocycle is not a coboundary.  This is the finite
full-shift instance of the periodic-data firewall familiar from Livšic theory
\citep{Livsic1972,Kalinin2011}.

\subsection{Honest countable domain}

The compatible countable atom limit is used only where its finite-fiber
weighted adjacency series converges absolutely.  For
\(\sigma_0=\Ree s>1\),
\[
\sum_{\varnothing\ne S\subset_{\rm fin}\mathbb P}|x_S|
=\prod_p(1+p^{-\sigma_0})-1<\infty.
\tag{3.8}\label{3.8}
\]
Consequently the series defining the \(2\times2\) fiber matrix converges in
operator norm, locally uniformly on \(\Ree s>1\), and defines a holomorphic
matrix family there.  We do not replace this finite-fiber quotient with an
unproved nuclear Ruelle operator on a larger countable-shift Banach space.
